\documentclass[11pt]{article}
\usepackage[margin=1in]{geometry}
\usepackage{amsmath, amssymb}
\usepackage{array}

\title{Matrix Equations and Inverses Cheat Sheet}
\author{}
\date{}

\begin{document}
\maketitle

\section*{9. Transformation Matrices: Graph the Image}

A transformation matrix changes a point or vector.

\[
\begin{bmatrix}
a & b\\
c & d
\end{bmatrix}
\begin{bmatrix}
x\\
y
\end{bmatrix}
=
\begin{bmatrix}
ax+by\\
cx+dy
\end{bmatrix}
\]

The result is the new point:

\[
(x,y)\rightarrow (x',y')
\]

Example:

\[
\begin{bmatrix}
2 & 0\\
0 & 2
\end{bmatrix}
\begin{bmatrix}
3\\
4
\end{bmatrix}
=
\begin{bmatrix}
2(3)+0(4)\\
0(3)+2(4)
\end{bmatrix}
=
\begin{bmatrix}
6\\
8
\end{bmatrix}
\]

So:

\[
(3,4)\rightarrow(6,8)
\]

\section*{10. Determinant of a Matrix}

For a \(2\times2\) matrix:

\[
A=
\begin{bmatrix}
a & b\\
c & d
\end{bmatrix}
\]

the determinant is:

\[
\det(A)=ad-bc
\]

Mental phrase:

\[
\text{diagonal down minus diagonal up}
\]

Example:

\[
A=
\begin{bmatrix}
3 & 5\\
2 & 7
\end{bmatrix}
\]

\[
\det(A)=3(7)-5(2)=21-10=11
\]

\section*{11. Is a Matrix Invertible?}

A matrix is invertible if its determinant is not zero:

\[
A \text{ is invertible if } \det(A)\ne 0
\]

\[
A \text{ is not invertible if } \det(A)=0
\]

Example:

\[
A=
\begin{bmatrix}
1 & 2\\
3 & 6
\end{bmatrix}
\]

\[
\det(A)=1(6)-2(3)=6-6=0
\]

So \(A\) is not invertible.

\section*{12. Inverse of a \(2\times2\) Matrix}

For:

\[
A=
\begin{bmatrix}
a & b\\
c & d
\end{bmatrix}
\]

the inverse is:

\[
A^{-1}
=
\frac{1}{ad-bc}
\begin{bmatrix}
d & -b\\
-c & a
\end{bmatrix}
\]

Steps:

\begin{enumerate}
  \item Find the determinant: \(ad-bc\).
  \item Swap \(a\) and \(d\).
  \item Make \(b\) and \(c\) negative.
  \item Divide every entry by the determinant.
\end{enumerate}

Example:

\[
A=
\begin{bmatrix}
1 & 2\\
3 & 4
\end{bmatrix}
\]

\[
\det(A)=1(4)-2(3)=4-6=-2
\]

Swap the diagonal and negate the off-diagonal:

\[
\begin{bmatrix}
4 & -2\\
-3 & 1
\end{bmatrix}
\]

Multiply by \(\frac{1}{-2}\):

\[
A^{-1}
=
\frac{1}{-2}
\begin{bmatrix}
4 & -2\\
-3 & 1
\end{bmatrix}
\]

\[
A^{-1}
=
\begin{bmatrix}
-2 & 1\\
\frac{3}{2} & -\frac{1}{2}
\end{bmatrix}
\]

\section*{13. Inverse of a \(3\times3\) Matrix}

The idea is:

\[
[A\mid I]\rightarrow[I\mid A^{-1}]
\]

Attach the identity matrix to \(A\), then row-reduce until the left side becomes the identity matrix.

Example setup:

\[
\left[
\begin{array}{ccc|ccc}
1 & 0 & 2 & 1 & 0 & 0\\
0 & 1 & 3 & 0 & 1 & 0\\
0 & 0 & 1 & 0 & 0 & 1
\end{array}
\right]
\]

Goal:

\[
\left[
\begin{array}{ccc|ccc}
1 & 0 & 0 & * & * & *\\
0 & 1 & 0 & * & * & *\\
0 & 0 & 1 & * & * & *
\end{array}
\right]
\]

The right side becomes \(A^{-1}\).

\section*{14. Identify Inverse Matrices}

The identity matrix is:

\[
I=
\begin{bmatrix}
1 & 0\\
0 & 1
\end{bmatrix}
\]

Two matrices are inverses if their product is the identity matrix:

\[
AB=I \Rightarrow A \text{ and } B \text{ are inverses}
\]

Example:

\[
A=
\begin{bmatrix}
2 & 0\\
0 & 3
\end{bmatrix}
\qquad
B=
\begin{bmatrix}
\frac{1}{2} & 0\\
0 & \frac{1}{3}
\end{bmatrix}
\]

\[
AB=
\begin{bmatrix}
2\cdot\frac{1}{2} & 0\\
0 & 3\cdot\frac{1}{3}
\end{bmatrix}
=
\begin{bmatrix}
1 & 0\\
0 & 1
\end{bmatrix}
\]

So \(A\) and \(B\) are inverses.

\section*{15. Solve Matrix Equations Using Inverses}

If:

\[
AX=B
\]

multiply by \(A^{-1}\) on the left:

\[
A^{-1}AX=A^{-1}B
\]

\[
IX=A^{-1}B
\]

\[
X=A^{-1}B
\]

If:

\[
XA=B
\]

multiply by \(A^{-1}\) on the right:

\[
XAA^{-1}=BA^{-1}
\]

\[
XI=BA^{-1}
\]

\[
X=BA^{-1}
\]

Important:

\[
\text{Matrix order matters.}
\]

\[
AX=B \Rightarrow X=A^{-1}B
\]

\[
XA=B \Rightarrow X=BA^{-1}
\]

\section*{Quick Checklist}

\[
\det
\begin{bmatrix}
a & b\\
c & d
\end{bmatrix}
=ad-bc
\]

\[
\det(A)\ne0 \Rightarrow A \text{ is invertible}
\]

\[
\det(A)=0 \Rightarrow A \text{ is not invertible}
\]

\[
\begin{bmatrix}
a & b\\
c & d
\end{bmatrix}^{-1}
=
\frac{1}{ad-bc}
\begin{bmatrix}
d & -b\\
-c & a
\end{bmatrix}
\]

\[
[A\mid I]\rightarrow[I\mid A^{-1}]
\]

\[
AB=I \Rightarrow A \text{ and } B \text{ are inverses}
\]

\[
AX=B \Rightarrow X=A^{-1}B
\]

\[
XA=B \Rightarrow X=BA^{-1}
\]

\end{document}
